Synthetic generators of equity-return time series support stress testing, regulatory backtesting, scenario design, and the training of machine-learning models on data that cannot be released for privacy or licensing reasons~\cite{jordon2022synthetic, assefa2020generating, alswaidan2026hybrid}. A useful generator has to reproduce, at the same time, the standard stylized facts of daily returns: heavy-tailed marginals, negligible linear autocorrelation, and slow decay of the autocorrelation function (ACF) of absolute returns~\cite{mandelbrot1963variation, cont2001empirical}. Generalised autoregressive conditional heteroskedasticity (GARCH) models~\cite{engle1982autoregressive, bollerslev1986generalized} capture volatility clustering but not the multi-regime structure of the marginal; i.i.d.\ generators match the marginal but lose temporal persistence; deep generators~\cite{yoon2019timegan, wiese2020quantgan, rasul2021autoregressive} reproduce some stylized facts at the cost of interpretability and reproducibility. One negative result has shaped two decades of work on regime-switching models. \citet{ryden1998stylized} fitted Gaussian-emission hidden Markov models (HMMs) at low state counts and concluded that the geometric sojourn-time distribution of a standard Markov chain makes the ACF decay too fast to match the observed slow decay; \citet{bulla2006stylized} recovered the ACF fit by replacing the Markov chain with a hidden semi-Markov chain, at the cost of much greater complexity. The Ryd\'{e}n et al.\ failure combines two constraints: a temporal one (the absolute-return ACF must allow slow-decay modes) and a distributional one (the marginal mixture must carry enough components to match observed kurtosis). A standard closed-form identity for the absolute-return ACF, written as a sum over the non-unit eigenvalues of the transition matrix~\cite{hamilton1994time, krolzig1997markov, timmermann2000moments}, bounds the number of decay modes by the rank of the centred transition matrix. We ask whether that bound actually matters on equity-return data once enough latent states are allowed.

In this study, we built three things on top of the continuous HMM (CHMM): a unified expectation-conditional-maximisation (ECM) framework for the four emission families, in which only the M-step changes between variants; a regime-conditional Value-at-Risk (VaR) that averages the predictive density over the one-step-ahead state forecast; and a Student-$t$ copula that ties the per-asset CHMM marginals into a joint generator through Sklar's theorem. We tested the model four ways: a six-fold rolling-origin walk-forward on the SPY decade, a sector-balanced 30-ticker panel across ten Global Industry Classification Standard (GICS) sectors, a Center for Research in Security Prices (CRSP) cross-decade transfer, and a six-asset US-equity basket, which together guard against single-window over-fitting. The four-family panel reproduced all three symmetric Cont stylized facts at the chosen $K^\star$, and heavy-tailed emissions narrowed the kurtosis gap relative to the Gaussian baseline. The rank bound from the bilinear ACF identity did not bind at the cross-ticker median, though it still mattered for a few tickers that introduced new regimes. The regime-conditional VaR passed the Christoffersen joint conditional-coverage test cleanly on the main window and on most walk-forward folds, and the Student-$t$ copula matched off-diagonal correlations far better than the single-index alternative. We read this to mean that the Ryd\'{e}n low-$K$ failure is a limit on the marginal mixture, not on the chain's temporal structure, once initialisation is fixed and the emissions can carry heavy tails. The i.i.d.\ bootstrap and a maximum-likelihood hidden semi-Markov model (HSMM) beat the CHMM on raw single-window distributional fit, but neither supports the regime-conditional risk forecasting, multi-asset copula composition, or parametric synthetic-data privacy the CHMM is built for. The scope is daily US equities: a non-equity stress test on the GLD exchange-traded fund (ETF) breaks down under static fitting, and we recommend periodic refit within that scope.
